\documentclass{tufte-handout}

\usepackage{amsmath}
\usepackage{booktabs}
\usepackage{graphicx}

\title{{{TITLE}}}
\author{Your Name}
\date{\today}

\begin{document}

\maketitle

\begin{abstract}
One paragraph: the problem, what you did, what you found, and why it matters. Keep it to about 150 words and avoid citations here.
\end{abstract}

\section{Introduction}

Tufte-style handouts keep the main text narrow and put notes, small figures
and references in the wide margin.\sidenote{Like this side note, numbered
like a footnote but placed next to the text it explains.} Citations appear in
the margin too.\cite{mcmahan2017communication}

\marginnote{A margin note has no number.}
Use them for asides that would interrupt the flow of the argument.

\section{Method}

The loss we minimise is
\[
  \mathcal{L}(\theta) = -\sum_{i=1}^{N} \log p_\theta(y_i \mid x_i).
\]

\begin{table}
  \centering
  \begin{tabular}{lrr}
    \toprule
    Method & Accuracy & Rounds \\
    \midrule
    Baseline & 81.2 & 500 \\
    Ours & 85.3 & 180 \\
    \bottomrule
  \end{tabular}
  \caption{A full-width table sits in the text column.}
\end{table}

\bibliography{references}
\bibliographystyle{plainnat}

\end{document}
